\section{LoRA}

Low-Rank Adaptation (LoRA) is a parameter-efficient\footnote{Parameter-efficient methods train relatively few task-specific parameters rather than updating every pre-trained model parameter.} fine-tuning\footnote{Fine-tuning adapts a pre-trained model to a downstream task using task-specific data.} method. It freezes pre-trained weights and learns low-rank updates for selected weight matrices. These notes follow \href{https://arxiv.org/abs/2106.09685}{Hu et al., \emph{LoRA: Low-Rank Adaptation of Large Language Models}}.

\subsection{Introduction}

\subsubsection{Problem}

Full fine-tuning updates all model parameters. For a large model such as GPT-3 (175 billion parameters), this makes training expensive and requires storing a separate full-size checkpoint for each task. Other parameter-efficient methods reduce task-specific storage, but some add inference latency or use part of the available context window.

\subsubsection{Motivation}

Li et al. define intrinsic dimension as the codimension of the solution set: if the full parameter space has dimension $D$ and the solution set has dimension $s$, then $D=d_{\mathrm{int}}+s$. Their random-subspace formulation keeps the initialization $\theta_0^{(D)}\in\mathbb R^D$ and projection $P\in\mathbb R^{D\times d}$ fixed while optimizing lower-dimensional coordinates $\theta^{(d)}\in\mathbb R^d$. At the final solution,
\[
	\boxed{\theta_*^{(D)}=\theta_0^{(D)}+P\theta_*^{(d)}}.
\]
Figure~\ref{fig:intrinsic_dimension} illustrates the case $D=3$, $s=1$, and $d=d_{\mathrm{int}}=2$. This intersection is a dimension-counting illustration: Li et al. establish almost-sure intersection at the codimension threshold for hyperplane solution sets, while arbitrary solution-set topologies need not intersect.

% ==========================================================
% Figure 1: Intrinsic dimension
% ==========================================================

\begin{figure}[htbp]
	\centering

	\begin{tikzpicture}[
			x=1cm,
			y=1cm,
			>=Stealth,
			font=\small,
			axis/.style={->,semithick},
			update/.style={->,semithick},
			solution/.style={
					line width=1.2pt,
					dash pattern=on .5pt off 3pt,
					line cap=round
				},
			plane/.style={
					fill=gray!25,
					fill opacity=1,
					draw=black!70,
					draw opacity=1,
					line width=.8pt
				},
			point/.style={
					circle,
					fill=black,
					inner sep=1.6pt
				}
		]

		% ------------------------------------------------------
		% (a) Full parameter space
		% ------------------------------------------------------
		\begin{scope}[shift={(0,0)}]

			% Negative theta_1 direction
			\draw[semithick] (-1.25,0) -- (0,0);

			% Positive theta_1 direction
			\draw[axis] (0,0) -- (2.15,0)
			node[right] {$\theta_1$};

			% Negative theta_2 direction
			\draw[semithick] (0,-.5) -- (0,0);

			% Positive theta_2 direction
			\draw[axis] (0,0) -- (0,2.9)
			node[above] {$\theta_2$};

			% Negative theta_3 direction
			\draw[semithick] (1.3,.86) -- (0,0);

			% Positive theta_3 direction
			\draw[axis] (0,0) -- (-1.3,-.86)
			node[below left] {$\theta_3$};

			% Initialization and final solution
			\coordinate (theta0L) at (.78,.65);
			\coordinate (thetaStarL) at (-.58,1.45);

			% Optimization displacement
			\draw[update] (theta0L) -- (thetaStarL);

			% Points
			\node[point] at (theta0L) {};
			\node[point] at (thetaStarL) {};

			% Labels
			\node[below right=2pt] at (theta0L)
			{$\theta_0^{(D)}$};

			\node[above left=2pt] at (thetaStarL)
			{$\theta_*^{(D)}$};

			% Panel heading
			\node[font=\bfseries] at (.35,3.65)
			{(a) Full parameter space};

			\node at (.4,-1.6) {$D=3$};

		\end{scope}


		% ------------------------------------------------------
		% (b) Affine optimization subspace
		% ------------------------------------------------------
		\begin{scope}[shift={(7.2,0)}]

			% --------------------------------------------------
			% Coordinate axes
			% --------------------------------------------------

			% Negative theta_1 direction
			\draw[semithick] (-1.25,0) -- (0,0);

			% Positive theta_1 direction
			\draw[axis] (0,0) -- (2.15,0)
			node[right] {$\theta_1$};

			% Negative theta_2 direction
			\draw[semithick] (0,-.5) -- (0,0);

			% Positive theta_2 direction
			\draw[axis] (0,0) -- (0,2.9)
			node[above right] {$\theta_2$};

			% Negative theta_3 direction
			\draw[semithick] (1.3,.86) -- (0,0);

			% Positive theta_3 direction
			\draw[axis] (0,0) -- (-1.3,-.86)
			node[below left] {$\theta_3$};

			% --------------------------------------------------
			% Opaque two-dimensional affine plane
			% --------------------------------------------------
			\coordinate (A) at (-.95,.05);
			\coordinate (B) at (2.25,.55);
			\coordinate (C) at (.70,2.35);
			\coordinate (E) at (-2.50,1.85);

			\filldraw[plane]
			(A) -- (B) -- (C) -- (E) -- cycle;

			% --------------------------------------------------
			% Initialization and final solution
			% --------------------------------------------------
			\coordinate (theta0R) at (.78,.65);
			\coordinate (thetaStarR) at (-.58,1.45);

			% Entry point of the solution line into the plane
			\coordinate (solutionEntry) at (-1.00687,.11603);

			% --------------------------------------------------
			% One-dimensional solution set
			% --------------------------------------------------

			% Solid portion below the surface
			\draw[semithick]
			(-1.156,-.35) -- (solutionEntry);

			% Sparse dotted portion behind the surface,
			% ending exactly at the final solution
			\draw[solution]
			(solutionEntry) -- (thetaStarR);

			% Solid portion above the final solution
			\draw[semithick]
			(thetaStarR) -- (-.116,2.9);

			% --------------------------------------------------
			% Optimization displacement inside the plane
			% --------------------------------------------------
			\draw[update]
			(theta0R) -- (thetaStarR)
			node[pos=.48,above right,
				fill=white,inner sep=2pt]
			{$P\theta_*^{(d)}$};

			% Points
			\node[point] at (theta0R) {};
			\node[point] at (thetaStarR) {};

			% Point labels
			\node[below right=2pt] at (theta0R)
			{$\theta_0^{(D)}$};

			\node[above left=2pt] at (thetaStarR)
			{$\theta_*^{(D)}$};

			% Panel heading
			\node[font=\bfseries] at (.35,3.65)
			{(b) Affine optimization subspace};

			% Parameterization
			\node[align=center] at (.35,-1.45) {
				$\theta^{(D)}=\theta_0^{(D)}+P\theta^{(d)}$
			};

			% Dimension relationship
			\node[align=center] at (.35,-2.15) {
				$D=d_{\mathrm{int}}+s=2+1=3$
			};

		\end{scope}

	\end{tikzpicture}

	\caption{
		Geometric illustration of intrinsic dimensionality,
		adapted from Li et al.~(2018).
		The opaque shaded affine plane has dimension $d=2$
		in the full parameter space $\mathbb R^3$.
		The one-dimensional solution set $\mathcal S$ is
		represented by a solid line outside the plane and
		a sparse dotted line inside it, intersecting the
		plane at $\theta_*^{(D)}$.
		In this example, $d_{\mathrm{int}}=D-s=2$.
		The intersection illustrates dimension counting
		and does not imply that every arbitrary solution set
		intersects every subspace of dimension $D-s$.
	}
	\label{fig:intrinsic_dimension}

\end{figure}

The low-dimensional fine-tuning results motivate LoRA, but do not establish that weight updates have low matrix rank. Figure~\ref{fig:lora_motivation} summarizes the progression from low-rank observations to LoRA's hypothesis and parameterization.


% ==========================================================
% Figure 2: Motivation for LoRA
% ==========================================================

\begin{figure}[htbp]
	\centering

	% IMPORTANT: Do not insert blank lines inside this
	% tikzpicture option list.
	\begin{tikzpicture}[
			>=Stealth,
			font=\small,
			node distance=5mm,
			box/.style={
					draw=black!70,
					fill=gray!8,
					rounded corners=2pt,
					text width=12.2cm,
					align=left,
					inner sep=7pt
				},
			hypothesis/.style={
					box,
					dashed,
					fill=gray!3
				},
			method/.style={
					box,
					fill=gray!17
				},
			arrow/.style={
					->,
					semithick
				}
		]

		% --------------------------------------------------
		% 1. Low-rank structure
		% --------------------------------------------------
		\node[box] (n1) {
			\textbf{1. Low-rank structure in machine learning}

			Many learning problems exhibit intrinsic low-rank
			structure. Overparameterized neural networks can
			also develop low-rank properties after training,
			and previous work has explicitly exploited these
			properties through low-rank constraints.

			\hfill{\footnotesize
				\href{https://arxiv.org/abs/1906.05392}
				{Oymak et al., 2019};
				\href{https://arxiv.org/abs/2106.09685}
				{Hu et al., 2021, Sec.~6}
			}
		};


		% --------------------------------------------------
		% 2. Intrinsic dimensionality of fine-tuning
		% --------------------------------------------------
		\node[box,below=of n1] (n2) {
			\textbf{2. Low intrinsic dimensionality of fine-tuning}

			Pretrained language models can be fine-tuned
			effectively through low-dimensional
			reparameterizations of their parameter space.

			\hfill{\footnotesize
				\href{https://aclanthology.org/2021.acl-long.568/}
				{Aghajanyan et al., 2021}
			}
		};


		% --------------------------------------------------
		% 3. LoRA hypothesis
		% --------------------------------------------------
		\node[hypothesis,below=of n2] (n3) {
			\textbf{3. LoRA hypothesis: low-rank weight updates}

			Do downstream tasks require only low-rank
			changes to pretrained weight matrices?

			\[
				\operatorname{rank}(\Delta W)
				\ll \min(m,n).
			\]

			Low-dimensional fine-tuning motivates this
			hypothesis, but does not imply it:
			intrinsic dimension and matrix rank are distinct
			quantities.

			\hfill{\footnotesize
				\href{https://arxiv.org/abs/2106.09685}
				{Hu et al., 2021; ICLR 2022}
			}
		};


		% --------------------------------------------------
		% 4. LoRA method and consequences
		% --------------------------------------------------
		\node[method,below=of n3] (n4) {
			\textbf{4. LoRA: low-rank parameterization}

			Freeze the pretrained matrix and optimize only
			its low-rank update:

			\[
				W=W_0+\Delta W=W_0+BA,
			\]

			\[
				W_0\in\mathbb R^{m\times n},
				\qquad
				B\in\mathbb R^{m\times r},
				\qquad
				A\in\mathbb R^{r\times n},
				\qquad
				r\ll\min(m,n).
			\]

			Fewer trainable parameters, reduced optimizer-state
			memory, task-specific adapters, and the ability to
			merge $BA$ into $W_0$ for inference.

			\hfill{\footnotesize
				\href{https://arxiv.org/abs/2106.09685}
				{Hu et al., 2021; ICLR 2022}
			}
		};


		% --------------------------------------------------
		% Connections
		% --------------------------------------------------
		\draw[arrow] (n1.south) -- (n2.north);

		% Dashed transition: motivation, not implication
		\draw[arrow,dashed] (n2.south) -- (n3.north);

		\draw[arrow] (n3.south) -- (n4.north);

	\end{tikzpicture}

	\caption{
		Conceptual progression motivating Low-Rank Adaptation.
		Observations of low-rank structure in machine learning
		and low intrinsic dimensionality in language-model
		fine-tuning motivate the hypothesis that task-specific
		weight updates have low intrinsic rank.
		The dashed transition indicates a hypothesis rather
		than a mathematical implication.
		LoRA implements this hypothesis through trainable
		low-rank factors while freezing pretrained weights.
	}
	\label{fig:lora_motivation}
\end{figure}

\subsubsection{Solution}

Motivated by evidence that model adaptation can take place in a low-dimensional subspace, the LoRA authors hypothesize that the required weight updates have low intrinsic rank.\footnote{Here, low intrinsic rank means that the task-specific change to a weight matrix can be represented effectively by a product of two matrices with a small shared dimension.} LoRA therefore represents the update to selected weight matrices with two trainable, low-rank matrices while keeping the original weights fixed.

This gives LoRA several practical advantages:
\begin{itemize}
	\item A single pre-trained model can be shared across tasks; each task needs only its own small LoRA matrices.
	\item Training uses fewer trainable parameters and less memory for gradients and optimizer states than full fine-tuning.
	\item The update can be merged into the original weights for inference, so it adds no inference latency.\footnote{Inference latency is the time a model takes to produce an output after receiving an input.}
	\item LoRA can be combined with other adaptation methods.
\end{itemize}

\subsection{Formulation}

\subsubsection{Problem statement}

Consider a pre-trained probabilistic model with parameters
\(\Phi_0\), which defines a conditional probability distribution \(P_{\Phi}(y\mid x)\).
For a downstream task, let \(\mathcal{D} = \{(x_i,y_i)\}_{i=1}^{N}\) be context-target pairs, where \(x_i\) and \(y_i\) are token sequences.
Full fine-tuning updates \(\Phi_0\) to maximize the conditional log-likelihood:
\[
	\min_{\Phi}
	- \sum_{(x,y)\in\mathcal{D}}
	\log P_{\Phi}(y \mid x)
\]

This approach learns a full-size parameter change \(\Delta\Phi\) for every task, which is extremely computationally expensive for large models.
LoRA instead learns a much smaller set of parameters \(\Theta\), which defines updates \(\Delta\Phi(\Theta)\) to selected weights while the pre-trained parameters stay frozen.
Let the task loss be
\[
	\mathcal{L}_{\mathrm{task}} =
	- \sum_{(x,y)\in\mathcal{D}}
	\log P_{\Phi_0+\Delta\Phi(\Theta)}(y \mid x)
\]
where \(|\Theta|\ll|\Phi_0|\).
Then the objective becomes
\[
	\min_{\Theta} \mathcal{L}_{\mathrm{task}}
\]

\subsubsection{Method}

For a pre-trained weight matrix \(W_0\in\mathbb{R}^{d\times k}\), LoRA constrains its update to \(\Delta W=AB\), where \(A\in\mathbb{R}^{d\times r}\), \(B\in\mathbb{R}^{r\times k}\), and \(r\ll\min(d,k)\). With input \(x\in\mathbb{R}^{k}\), the adapted layer computes
\[
	h=W_0x+\frac{\alpha}{r}ABx
\]
where \(\alpha\) is a scaling hyperparameter.
The factor \(\alpha/r\) reduces the need to retune other hyperparameters when the rank \(r\) changes.
Only \(A\) and \(B\) are trained; \(W_0\) remains frozen.

For this matrix, LoRA trains \(r(d+k)\) parameters instead of the \(dk\) weights that full fine-tuning would update. This reduces the task-specific checkpoint size and the gradient and optimizer-state memory for the adapted matrix. The base weights must still be stored and used.

\subsubsection{Initialization and deployment}

The original LoRA method initializes \(A\) with random Gaussian values and \(B\) with zeros. Consequently, \(AB=0\) at the start of training, so the adapted layer initially computes exactly \(W_0x\). Training then changes \(A\) and \(B\) while leaving \(W_0\) fixed.

After training, the update can be merged into the weight matrix:
\[
	W_{\mathrm{merged}}=W_0+\frac{\alpha}{r}AB
\]
Inference with \(W_{\mathrm{merged}}x\) uses the same matrix multiplication as the original layer and incurs no additional latency from a separate LoRA path. For multiple tasks, one can retain the shared base model and store a small adapter for each task, merging the chosen adapter when needed. Deploying an adapter separately still requires access to its matching base model.
